\chapter{K-Means Variants, BIC Criterion, Mixture Models, and EM}
\label{chap:dm-kmeans-variants-bic-em}

In this chapter, we explore the theoretical and algorithmic evolution that overcomes the intrinsic limitations of basic standard $K$-Means (Lloyd's algorithm), introducing advanced divisive partitioning techniques, formal information-theoretic criteria for automatic model complexity selection, and probabilistic generative models founded on the \textit{Expectation-Maximization} (EM) paradigm.
The presentation rigorously covers:
\begin{enumerate}[noitemsep]
    \item The \textbf{structural and geometric limitations of standard $K$-Means}: extreme sensitivity to initial centroid seeding, the rigidity of deterministic assignment (\textit{hard clustering}), the assumption of hyperspherical clusters with isotropic and homogeneous variances, and the requirement to specify the hyperparameter $K$ a priori;
    \item The \textbf{Bisecting $K$-Means} algorithm: the transition to the divisive hierarchical (\textit{top-down}) paradigm, decomposing the global partitioning problem into a sequence of local two-way bisections ($K=2$), multi-trial execution with minimization of intra-cluster sum of squared errors ($\mathrm{SSE}$), split selection heuristics, boundary/singleton cases, and post-processing global Lloyd refinement;
    \item The \textbf{$X$-Means} algorithm (Pelleg and Moore~\cite{pelleg2000}): modeling clustering as a statistical model selection problem over the interval $[K_{\min}, K_{\max}]$, local bisection testing guided by principal axis perturbations, local $2$-means optimization, and information gain validation via $\Delta\mathrm{BIC}$;
    \item The \textbf{Bayesian Information Criterion (BIC)} (Schwarz~\cite{schwarz1978}): asymptotic derivation from the marginal likelihood via Laplace's approximation, the analytical trade-off between log-likelihood and complexity penalty (Occam's razor), and precise counting of free parameters under spherical, diagonal, and full covariance matrices;
    \item The \textbf{Gaussian Mixture Model (GMM)} framework: extending clustering to the probabilistic domain (\textit{soft clustering}), generative formulation using latent categorical variables, incomplete vs. complete data log-likelihood, and Maximum Likelihood Estimation (MLE) vs. Maximum A Posteriori (MAP);
    \item The foundational \textbf{Expectation-Maximization (EM)} algorithm (Dempster, Laird, and Rubin~\cite{dempster1977}): Jensen's inequality, the variational evidence lower bound, the auxiliary $\mathcal{Q}$-function, and the cyclic alternation between the E-step (analytical posterior responsibilities $\gamma_{ik}$) and the M-step (closed-form weighted updates for mixing proportions $\pi_k$, mean vectors $\boldsymbol{\mu}_k$, and covariance matrices $\boldsymbol{\Sigma}_k$);
    \item A \textbf{concrete analytical case study on a two-Gaussian mixture}: deconvolution of overlapping demographic populations (e.g., human heights), comparison between theoretical continuous PDF and empirical 20,000-sample histogram, critical saddle point analysis ($x=0$), and soft uncertainty assignment;
    \item A \textbf{comprehensive architectural comparison between $K$-Means and EM}: geometric and energetic taxonomy, operational strengths, computational vulnerabilities, covariance singularities, and high-dimensional regularization.
\end{enumerate}

% ====================================================================
% SECTION 9.1: SCIENTIFIC CONTEXT AND LIMITS OF STANDARD K-MEANS
% ====================================================================
\section{Scientific Context and Limitations of Standard \texorpdfstring{$K$}{K}-Means}
\label{sec:dm-kmeans-limits}

Standard $K$-Means (Lloyd's algorithm), formalized in Chapter~\ref{chap:dm-kmeans-hierarchical} (\autoref{sec:dm-kmeans-algorithm-geometry}), operates by locating $K$ centroids $\mathbf{C} = \{\mathbf{c}_1, \dots, \mathbf{c}_K\}$ that minimize the global intra-cluster Sum of Squared Errors $\mathrm{SSE}$ (\autoref{eq:dm-sse-definition-en}). While computationally linear per iteration ($\mathcal{O}(N \cdot K \cdot d)$), the structural analysis in \autoref{sec:dm-kmeans-limitations-isodata} highlighted three primary vulnerabilities that constrain its utility in practical applications:

\begin{enumerate}
    \item \textbf{Extreme sensitivity to initialization}: Continuous $\mathrm{SSE}$ minimization is non-convex and $\mathcal{NP}$-hard. Lloyd's coordinate descent guarantees convergence only to a local minimum: seeds placed sub-optimally or within the same natural dense region trap the algorithm in fragmented, high-error states;
    \item \textbf{Requirement of predetermining parameter $K$}: Standard $K$-Means lacks internal model selection. Because $\mathrm{SSE}$ decreases monotonically with $K$, it cannot independently guide the selection of the natural cluster count without complexity penalties;
    \item \textbf{Geometric rigidity and deterministic assignment (\textit{hard clustering})}: $K$-Means forces every instance into a single, mutually exclusive cluster, carving out linear Voronoi cells. This formulation implicitly assumes clusters are hyperspherical, convex, equally sized, and endowed with identical isotropic variances.
\end{enumerate}

\begin{figure}[htbp]
\centering
\begin{tikzpicture}[
    boxStyle/.style={rectangle, rounded corners=6pt, draw=navy!85!black, fill=navy!6, thick, font=\footnotesize, text width=3.8cm, inner sep=6pt, align=left, minimum height=3.5cm},
    headerStyle/.style={rectangle, rounded corners=6pt, draw=purple!85!black, fill=purple!12, very thick, font=\small, align=center, text width=12.8cm, inner sep=7pt},
    arrowStyle/.style={-Stealth, thick, draw=navy!85!black}
]
    \node[headerStyle] (top) at (0, 4.6) {
        \textbf{\large LIMITATIONS OF $K$-MEANS AND REMEDIES}\\[3pt]
        \footnotesize Transitioning from rigid geometric heuristics to adaptive and probabilistic paradigms
    };

    \node[boxStyle, anchor=north] (b1) at (-4.5, 1.6) {
        \centering\textbf{Seed Sensitivity}\\[2pt]
        \centering\scriptsize (Sub-optimal $\mathrm{SSE}$ local minima)\\[5pt]
        \raggedright
        $\bullet$ Poor convergence risk\\[2pt]
        $\bullet$ Remedy: \textbf{Bisecting K-Means}\\[2pt]
        $\bullet$ Divisive hierarchical tree\\[2pt]
        $\bullet$ Multi-trial bisections
    };

    \node[boxStyle, anchor=north] (b2) at (0, 1.6) {
        \centering\textbf{Preset $K$ Requirement}\\[2pt]
        \centering\scriptsize (No intrinsic model selection)\\[5pt]
        \raggedright
        $\bullet$ No objective criteria\\[2pt]
        $\bullet$ Remedy: \textbf{X-means and BIC}\\[2pt]
        $\bullet$ Schwarz complexity penalty\\[2pt]
        $\bullet$ Local statistical splitting
    };

    \node[boxStyle, anchor=north] (b3) at (4.5, 1.6) {
        \centering\textbf{Rigid Voronoi Hard Cut}\\[2pt]
        \centering\scriptsize (Spherical isotropic assumption)\\[5pt]
        \raggedright
        $\bullet$ Zero boundary uncertainty\\[2pt]
        $\bullet$ Remedy: \textbf{GMM and EM}\\[2pt]
        $\bullet$ Probabilistic \textit{soft} clustering\\[2pt]
        $\bullet$ Heterogeneous covariance $\boldsymbol{\Sigma}_k$
    };

    % Branching connectors with ample spacing and rounded corners
    \draw[thick, draw=navy!85!black] (top.south) -- (0, 3.0);
    \draw[arrowStyle, rounded corners=4pt] (0, 3.0) -| (b1.north);
    \draw[arrowStyle] (0, 3.0) -- (b2.north);
    \draw[arrowStyle, rounded corners=4pt] (0, 3.0) -| (b3.north);
\end{tikzpicture}
\caption{Conceptual overview of the foundational limitations of standard $K$-Means and the three complementary algorithmic strategies designed to address them.}
\label{fig:dm-kmeans-limits-overview}
\end{figure}

To overcome these structural hurdles, data mining research has developed three complementary methodologies illustrated in Figure~\ref{fig:dm-kmeans-limits-overview}:
\begin{itemize}
    \item \textbf{Bisecting $K$-Means}: a top-down divisive hierarchical framework that drastically mitigates initialization fragility by converting global clustering into a sequence of tractable binary sub-problems ($K=2$);
    \item \textbf{$X$-Means and the BIC Criterion}: incorporating formal information-theoretic regularization to autonomously discover the intrinsic number of clusters;
    \item \textbf{Gaussian Mixture Models and Expectation-Maximization}: replacing rigid geometric tessellations with a generative probabilistic model featuring continuous degree-of-belonging values (\textit{posterior responsibilities}).
\end{itemize}

% ====================================================================
% SECTION 9.2: BISECTING K-MEANS
% ====================================================================
\section{The Bisecting \texorpdfstring{$K$}{K}-Means Algorithm}
\label{sec:dm-bisecting-kmeans}

The \textbf{Bisecting $K$-Means} algorithm (formalized as a divisive variant in~\cite{tan2018}) addresses the initialization problem by replacing simultaneous global seeding with a top-down hierarchical procedure founded on the \textit{divide-and-conquer} principle. Instead of positioning $K$ centroids simultaneously in high-dimensional space, the algorithm decomposes the challenge into $K-1$ sequential, independent two-way bisections ($K=2$).

\subsection{Mathematical Formulation of Error and Partition Strategy}

Given a subset of instances $\mathcal{C} \subseteq \mathcal{D}$ with cardinality $|\mathcal{C}| = n_{\mathcal{C}}$, the intra-cluster error $\mathrm{SSE}(\mathcal{C})$ is computed relative to its centroid $\mathbf{c} = \frac{1}{n_{\mathcal{C}}} \sum_{\mathbf{x} \in \mathcal{C}} \mathbf{x}$:
\begin{equation}
\label{eq:dm-sse-single-cluster}
\mathrm{SSE}(\mathcal{C}) = \sum_{\mathbf{x} \in \mathcal{C}} \|\mathbf{x} - \mathbf{c}\|_2^2
\end{equation}
If the current partition of the space is represented by an active list of clusters $\mathcal{L} = \{\mathcal{C}_1, \mathcal{C}_2, \dots, \mathcal{C}_m\}$, the total sum of squared errors across the configuration is:
\begin{equation}
\label{eq:dm-sse-total-list}
\mathrm{SSE}_{\text{total}}(\mathcal{L}) = \sum_{j=1}^m \mathrm{SSE}(\mathcal{C}_j)
\end{equation}
When a candidate cluster $\mathcal{C}^* \in \mathcal{L}$ is selected, extracted, and bisected into two disjoint child clusters $\mathcal{A}$ and $\mathcal{B}$ (such that $\mathcal{A} \cup \mathcal{B} = \mathcal{C}^*$ and $\mathcal{A} \cap \mathcal{B} = \emptyset$), the updated global sum of squared errors becomes:
\begin{equation}
\label{eq:dm-sse-split-delta}
\mathrm{SSE}_{\text{total}}(\mathcal{L}') = \mathrm{SSE}_{\text{total}}(\mathcal{L}) - \mathrm{SSE}(\mathcal{C}^*) + \big( \mathrm{SSE}(\mathcal{A}) + \mathrm{SSE}(\mathcal{B}) \big)
\end{equation}
By the Huygens-Steiner theorem of variance decomposition, the dispersion about two independent centroids is strictly no greater than the dispersion about a single pooled centroid: $\mathrm{SSE}(\mathcal{A}) + \mathrm{SSE}(\mathcal{B}) \le \mathrm{SSE}(\mathcal{C}^*)$. Consequently, every bisection step monotonically decreases the global residual error.

\subsection{Pseudocode and Procedural Analysis}

Algorithm~\ref{alg:dm-bisecting-kmeans} formalizes the procedural workflow of Bisecting $K$-Means.

\begin{algorithm}[htbp]
\DontPrintSemicolon
\caption{Bisecting $K$-Means Algorithm}
\label{alg:dm-bisecting-kmeans}
\KwInput{Dataset $\mathcal{D} = \{\mathbf{x}_1, \dots, \mathbf{x}_N\} \subset \mathbb{R}^d$, desired number of clusters $K$, number of bisection trials $T$}
\KwOutput{Final list of $K$ clusters $\mathcal{L} = \{\mathcal{C}_1, \dots, \mathcal{C}_K\}$}
Initialize cluster list $\mathcal{L} \leftarrow \{\mathcal{D}\}$ containing the entire dataset as a single root cluster\;
\While{$|\mathcal{L}| < K$}{
    Select and remove a candidate cluster $\mathcal{C}^*$ from $\mathcal{L}$ using a split heuristic (e.g., maximum $\mathrm{SSE}$ or maximum cardinality)\;
    $\mathcal{L} \leftarrow \mathcal{L} \setminus \{\mathcal{C}^*\}$\;
    $\text{best\_sse} \leftarrow +\infty$\;
    $(\mathcal{A}^*, \mathcal{B}^*) \leftarrow (\emptyset, \emptyset)$\;
    \For{$i \leftarrow 1$ \KwTo $T$}{
        Randomly initialize two centroids within $\mathcal{C}^*$\;
        Execute basic $K$-Means with $K=2$ on $\mathcal{C}^*$, yielding child clusters $(\mathcal{A}_i, \mathcal{B}_i)$\;
        $\text{current\_sse} \leftarrow \mathrm{SSE}(\mathcal{A}_i) + \mathrm{SSE}(\mathcal{B}_i)$\;
        \If{$\text{current\_sse} < \text{best\_sse}$}{
            $\text{best\_sse} \leftarrow \text{current\_sse}$\;
            $(\mathcal{A}^*, \mathcal{B}^*) \leftarrow (\mathcal{A}_i, \mathcal{B}_i)$\;
        }
    }
    Insert optimal child clusters into list: $\mathcal{L} \leftarrow \mathcal{L} \cup \{\mathcal{A}^*, \mathcal{B}^*\}$\;
}
\Return $\mathcal{L}$\;
\end{algorithm}

\subsection{Cluster Selection Heuristics and Multi-Trial Execution}

The algorithm provides two essential degrees of freedom:
\begin{enumerate}
    \item \textbf{Candidate Cluster Selection Strategy $\mathcal{C}^*$}:
    \begin{itemize}
        \item \textit{Maximum Individual SSE} ($\mathcal{C}^* = \arg\max_{\mathcal{C} \in \mathcal{L}} \mathrm{SSE}(\mathcal{C})$): The standard recommended strategy in literature. By splitting the cluster exhibiting the highest internal variance, one maximizes the local error reduction $\mathrm{SSE}(\mathcal{C}^*) - (\mathrm{SSE}(\mathcal{A}) + \mathrm{SSE}(\mathcal{B}))$, actively promoting variance equalization across clusters.
        \item \textit{Maximum Cardinality} ($\mathcal{C}^* = \arg\max_{\mathcal{C} \in \mathcal{L}} |\mathcal{C}|$): Selects the most populous group, particularly effective when data contains significant cluster size imbalances or oversized clusters.
    \end{itemize}
    \item \textbf{Role of Multiple Trials ($T$ Trials)}: Because basic $2$-means remains an iterative heuristic sensitive to initial seed placement within $\mathcal{C}^*$, performing $T$ independent trials (typically $T \in [5, 20]$) and retaining the partition minimizing $\mathrm{SSE}(\mathcal{A}) + \mathrm{SSE}(\mathcal{B})$ exponentially reduces the likelihood of adopting a sub-optimal split.
\end{enumerate}

\begin{figure}[htbp]
\centering
\providecolor{accentblue}{HTML}{1976D2}
\providecolor{purple}{HTML}{7B1FA2}
\begin{tikzpicture}[
    treeNode/.style={rectangle, rounded corners=5pt, draw=navy!85!black, fill=navy!8, thick, font=\footnotesize, align=center, inner sep=6pt, text width=3.6cm},
    leafNode/.style={rectangle, rounded corners=5pt, draw=darkgreen!85!black, fill=darkgreen!10, thick, font=\footnotesize, align=center, inner sep=6pt, text width=3.5cm},
    leafBlue/.style={rectangle, rounded corners=5pt, draw=accentblue!85!black, fill=accentblue!10, thick, font=\footnotesize, align=center, inner sep=5pt, text width=3.1cm},
    leafNavy/.style={rectangle, rounded corners=5pt, draw=navy!85!black, fill=navy!10, thick, font=\footnotesize, align=center, inner sep=5pt, text width=3.1cm},
    activeNode/.style={rectangle, rounded corners=5pt, draw=crimson!85!black, fill=crimson!10, thick, font=\footnotesize, align=center, inner sep=6pt, text width=3.5cm},
    badgeTrial/.style={rectangle, rounded corners=4pt, draw=purple!80!black, fill=purple!8, dashed, thick, font=\scriptsize, align=center, inner sep=5pt, text width=6.6cm},
    arrowTree/.style={-Stealth, thick, draw=navy!85!black, rounded corners=3pt}
]

    % Level 0: Root (centered at x = 0, y = 7.4)
    \node[treeNode] (root) at (0, 7.4) {
        \textbf{Root Cluster $\mathcal{D}$}\\
        \scriptsize $N$ samples $\cdot$ $\mathrm{SSE}_{\text{tot}}$
    };

    % Operation 1: Split 1 (centered at x = 0, y = 5.5)
    \node[badgeTrial] (split1) at (0, 5.5) {
        \textbf{Split 1}: $T$ trials with $K=2$ on $\mathcal{D}$\\[2pt]
        \tiny Min $[\mathrm{SSE}(\mathcal{C}_1) + \mathrm{SSE}(\mathcal{C}_2)]$
    };

    % Level 1: C1 and C2 (symmetric at x = -3.7 and x = +3.7, y = 2.8)
    \node[activeNode] (c1) at (-3.7, 2.8) {
        \textbf{Cluster $\mathcal{C}_1$ [Max SSE]}\\
        \scriptsize $\mathrm{SSE}(\mathcal{C}_1) > \mathrm{SSE}(\mathcal{C}_2)$\\[2pt]
        \textbf{\textcolor{crimson!90!black}{Selected for bisection}}
    };

    \node[leafNode] (c2) at (3.7, 2.8) {
        \textbf{Cluster $\mathcal{C}_2$ [Leaf 1]}\\
        \scriptsize Low $\mathrm{SSE}$ $\cdot$ compact\\[2pt]
        \textbf{\textcolor{darkgreen!90!black}{Preserved intact in $\mathcal{L}$}}
    };

    % Operation 2: Split 2 (centered under C1 at x = -3.7, y = 0.5)
    \node[badgeTrial, text width=6.0cm] (split2) at (-3.7, 0.5) {
        \textbf{Split 2}: $T$ trials with $K=2$ on $\mathcal{C}_1$\\[2pt]
        \tiny Min $[\mathrm{SSE}(\mathcal{C}_{1a}) + \mathrm{SSE}(\mathcal{C}_{1b})]$
    };

    % Level 2: Leaves C1a and C1b (centered around x = -3.7, y = -2.2)
    \node[leafBlue] (c1a) at (-5.5, -2.2) {
        \textbf{Cluster $\mathcal{C}_{1a}$ [Leaf 2]}\\[2pt]
        \scriptsize Optimal sub-cluster\\
        \scriptsize Minimal residual $\mathrm{SSE}$
    };

    \node[leafNavy] (c1b) at (-1.9, -2.2) {
        \textbf{Cluster $\mathcal{C}_{1b}$ [Leaf 3]}\\[2pt]
        \scriptsize Optimal sub-cluster\\
        \scriptsize Minimal residual $\mathrm{SSE}$
    };

    % C2 baseline state node (aligned at x = +3.7, y = -2.2)
    \node[leafNode] (c2leaf) at (3.7, -2.2) {
        \textbf{Cluster $\mathcal{C}_2$ [Leaf 1]}\\[2pt]
        \scriptsize Unchanged state\\
        \scriptsize No further split required
    };

    % --- Connectors ---
    % Root -> Split 1
    \draw[arrowTree] (root.south) -- (split1.north);

    % Split 1 -> C1 and C2 (orthogonal branching downwards)
    \draw[thick, draw=navy!85!black] (split1.south) -- (0, 4.3);
    \draw[arrowTree] (0, 4.3) -| (c1.north);
    \draw[arrowTree] (0, 4.3) -| (c2.north);

    % C1 -> Split 2
    \draw[arrowTree] (c1.south) -- (split2.north);

    % Split 2 -> C1a and C1b (orthogonal branching downwards)
    \draw[thick, draw=navy!85!black] (split2.south) -- (-3.7, -0.7);
    \draw[arrowTree] (-3.7, -0.7) -| (c1a.north);
    \draw[arrowTree] (-3.7, -0.7) -| (c1b.north);

    % C2 continuation dashed arrow
    \draw[dashed, draw=darkgreen!80!black, very thick, -Stealth] (c2.south) -- (c2leaf.north)
        node[midway, right=7pt, font=\scriptsize, text=darkgreen!85!black, align=left] {
            High homogeneity:\\
            branch complete
        };

    % Bottom Summary Banner (y = -3.9)
    \node[anchor=north, rectangle, rounded corners=5pt, draw=gray!40, fill=gray!4, font=\scriptsize, align=center, text width=13.0cm, inner sep=6pt] at (0, -3.9) {
        \textbf{Top-down divisive progression}: greedily select the cluster with maximum $\mathrm{SSE}$ from $\mathcal{D}$.\\
        Execute $T$ trials with $K=2$, choosing the split that minimizes $\mathrm{SSE}(\mathcal{A}) + \mathrm{SSE}(\mathcal{B})$.\\
        Homogeneous clusters remain leaves, prioritizing dispersed groups until reaching $K$ clusters.
    };
\end{tikzpicture}
\caption{Hierarchical binary tree representation of the divisive progression of Bisecting $K$-Means with multiple bisection trials and minimum $\mathrm{SSE}$ selection.}
\label{fig:dm-bisecting-kmeans-tree}
\end{figure}


\subsection{Computational Complexity and Boundary Cases}

Computational analysis reveals significant operational advantages:
\begin{itemize}
    \item \textbf{Per-Bisection Complexity}: Running $2$-means on a cluster of $n_k$ points for $I_2$ iterations requires time $\mathcal{O}(I_2 \cdot 2 \cdot n_k \cdot d)$. Because at any tree depth the sum of cluster sizes is $\sum n_k = N$, if the divisive tree is reasonably balanced the overall depth is $\mathcal{O}(\log K)$. The total cost amounts to $\mathcal{O}(T \cdot I_2 \cdot N \cdot d \cdot \log K)$, which is significantly faster than standard global $K$-Means when $K$ is large.
    \item \textbf{Degenerate Clusters and Singletons}: If a candidate cluster contains only one point ($|\mathcal{C}| = 1$) or has zero variance (identical points), it cannot be bisected. The algorithm skips such clusters and inspects the next candidate ordered by $\mathrm{SSE}$.
\end{itemize}

\begin{noteblock}{Global Lloyd Refinement}
Because partitioning decisions in Bisecting $K$-Means are made greedily and locally (points in $\mathcal{C}^*$ are separated without interacting with points in other existing clusters), decision boundaries may slightly diverge from global optimality. It is standard best practice to initialize a final run of standard $K$-Means using the $K$ centroids produced by Bisecting $K$-Means. This refinement typically converges in very few iterations (2--4 iterations), cleaning up boundary assignments without any risk of getting trapped in bad local minima.
\end{noteblock}

% ====================================================================
% SECTION 9.3: X-MEANS AND THE BIC CRITERION
% ====================================================================
\section{The \texorpdfstring{$X$}{X}-Means Algorithm and Model Selection}
\label{sec:dm-xmeans-algorithm}

The \textbf{$X$-Means} algorithm, introduced by Dan Pelleg and Andrew Moore~\cite{pelleg2000}, addresses the fundamental problem of choosing $K$ automatically and rigorously. Instead of asking the user for a rigid fixed value, $X$-Means requires only a bounding search range $[K_{\min}, K_{\max}]$, dynamically exploring model complexity via local statistical tests guided by the \textbf{Bayesian Information Criterion} (BIC).

\subsection{Procedural Flow and Local Split Decisions}

$X$-Means alternates between global clustering iterations and statistical split evaluation:
\begin{enumerate}
    \item \textbf{Initialization and Global Convergence}: The process initializes at $K = K_{\min}$ and executes conventional $K$-Means to convergence.
    \item \textbf{Centroid Perturbation along Principal Direction}: For each active cluster $\mathcal{C}_k$ with centroid $\mathbf{c}_k$, the algorithm evaluates whether $\mathcal{C}_k$ should remain a single entity (null hypothesis $\mathcal{M}_1$) or be divided into two child clusters (alternative hypothesis $\mathcal{M}_2$). Two candidate seeds $\mathbf{c}_k^{(1)}$ and $\mathbf{c}_k^{(2)}$ are generated by perturbing the centroid along direction $\mathbf{v}$ proportional to the local principal axis of variance:
    \begin{equation}
    \label{eq:dm-xmeans-perturbation}
    \mathbf{c}_k^{(1)} = \mathbf{c}_k - \alpha \mathbf{v}, \qquad \mathbf{c}_k^{(2)} = \mathbf{c}_k + \alpha \mathbf{v}
    \end{equation}
    where $\alpha = s \cdot \sqrt{\lambda_{\max}}$, with $\lambda_{\max}$ being the maximum eigenvalue of the local covariance and $s$ a scaling constant.
    \item \textbf{Local $2$-means Clustering}: A local instance of $2$-means is executed strictly on points belonging to $\mathcal{C}_k$, yielding refined child centroids $\hat{\mathbf{c}}_k^{(1)}$ and $\hat{\mathbf{c}}_k^{(2)}$.
    \item \textbf{BIC Model Evaluation}: The Bayesian Information Criterion is computed for the single-cluster hypothesis $\mathrm{BIC}(\mathcal{M}_1)$ and the two-cluster hypothesis $\mathrm{BIC}(\mathcal{M}_2)$. The information gain is evaluated as:
    \begin{equation}
    \label{eq:dm-delta-bic}
    \Delta\mathrm{BIC} = \mathrm{BIC}(\mathcal{M}_2) - \mathrm{BIC}(\mathcal{M}_1)
    \end{equation}
    If $\Delta\mathrm{BIC} > 0$, the gain in likelihood outweighs the complexity penalty of additional parameters: the split is accepted, and $\mathbf{c}_k$ is replaced by the pair $(\hat{\mathbf{c}}_k^{(1)}, \hat{\mathbf{c}}_k^{(2)})$. Otherwise, the split is rejected.
    \item \textbf{Termination Check}: If no clusters are split or if the total number of centroids reaches $K_{\max}$, the algorithm terminates. Otherwise, global $K$-Means is re-run with the expanded centroid set, and the cycle repeats.
\end{enumerate}

\begin{figure}[htbp]
\centering
\begin{tikzpicture}[
    panelBox/.style={rectangle, rounded corners=4.5pt, draw=navy!80!black, fill=navy!4, thick, font=\footnotesize, inner sep=6pt, text width=5.7cm, minimum height=3.3cm, align=left},
    tagBadge/.style={rectangle, rounded corners=3.5pt, draw=purple!85!black, fill=purple!12, thick, font=\scriptsize\bfseries, text=purple!90!black, inner sep=2.5pt}
]
    % Panel 1: Top-Left
    \node[panelBox] (p1) at (-3.6, 2.3) {
        \textbf{(a) Initial Voronoi Partition}\\[3pt]
        $\bullet$ Run $K$-Means with $K=3$.\\
        $\bullet$ Metric space partitioned into Voronoi cells $V_k$.\\
        $\bullet$ Points assigned to nearest centroid $\mathbf{c}_k$.
    };
    \node[tagBadge, anchor=south east] at ([xshift=-10pt, yshift=-1.5pt]p1.north east) {Stage 1};

    % Panel 2: Top-Right
    \node[panelBox] (p2) at (3.6, 2.3) {
        \textbf{(b) Centroid Perturbation}\\[3pt]
        $\bullet$ Generate 2 child seeds per cell.\\
        $\bullet$ Opposite shift along principal axis:\\
        $\quad \mathbf{c}_k^{(1)} = \mathbf{c}_k - \alpha \mathbf{v}$\\
        $\quad \mathbf{c}_k^{(2)} = \mathbf{c}_k + \alpha \mathbf{v}$\\
        $\bullet$ Prepare for local trial bisection.
    };
    \node[tagBadge, anchor=south east] at ([xshift=-10pt, yshift=-1.5pt]p2.north east) {Stage 2};

    % Panel 3: Bottom-Left
    \node[panelBox] (p3) at (-3.6, -2.1) {
        \textbf{(c) Local $2$-means Bisection}\\[3pt]
        $\bullet$ $2$-means run strictly on data in $V_k$.\\
        $\bullet$ Rapid optimization of centroid positions.\\
        $\bullet$ Zero interference between adjacent regions (parallelizable).
    };
    \node[tagBadge, anchor=south east] at ([xshift=-10pt, yshift=-1.5pt]p3.north east) {Stage 3};

    % Panel 4: Bottom-Right
    \node[panelBox] (p4) at (3.6, -2.1) {
        \textbf{(d) Local BIC Evaluation}\\[3pt]
        $\bullet$ Evaluate $\Delta\mathrm{BIC} = \mathrm{BIC}(\mathcal{M}_2) - \mathrm{BIC}(\mathcal{M}_1)$.\\
        $\bullet$ \textcolor{darkgreen}{\textbf{Split Accepted}}: if $\Delta\mathrm{BIC} > 0$ (upper cluster non-spherical).\\
        $\bullet$ \textcolor{crimson}{\textbf{Split Rejected}}: if $\Delta\mathrm{BIC} \le 0$ (lower clusters already compact).
    };
    \node[tagBadge, anchor=south east] at ([xshift=-10pt, yshift=-1.5pt]p4.north east) {Stage 4};

    % Central sequence arrows
    \draw[-Stealth, very thick, draw=navy!80!black] (p1.east) -- (p2.west);
    \draw[-Stealth, very thick, draw=navy!80!black] (p2.south) -- (p4.north);
    \draw[-Stealth, very thick, draw=navy!80!black] (p4.west) -- (p3.east);
    \draw[-Stealth, very thick, draw=navy!80!black, dashed] (p3.north) -- (p1.south) node[midway, left=4pt, font=\scriptsize\bfseries, text=navy!80!black] {Global Loop};
\end{tikzpicture}
\caption{Four-stage operational cycle of $X$-Means: initial Voronoi partition, directional seed perturbation, local bisection, and statistical validation via the information-theoretic $\Delta\mathrm{BIC}$ test.}
\label{fig:dm-xmeans-cycle}
\end{figure}

% ====================================================================
% SECTION 9.4: THE BAYESIAN INFORMATION CRITERION (BIC)
% ====================================================================
\section{The Bayesian Information Criterion (BIC)}
\label{sec:dm-bic-theory}

The Bayesian Information Criterion (BIC), formulated by Gideon Schwarz~\cite{schwarz1978}, is grounded in asymptotic Bayesian model selection theory. It provides a principled mathematical mechanism for trading off \textbf{goodness of fit} against \textbf{parametric model complexity} (\textit{parsimony or Occam's razor}).

\subsection{Asymptotic Derivation from Marginal Likelihood}

Consider an observed dataset $\mathcal{D} = \{\mathbf{x}_1, \dots, \mathbf{x}_R\}$ and a parametric model family $\mathcal{M}$ indexed by parameter vector $\boldsymbol{\theta} \in \mathbb{R}^p$. In Bayesian inference, the posterior probability of model $\mathcal{M}$ given $\mathcal{D}$ is given by Bayes' theorem:
\begin{equation}
\label{eq:dm-bayes-model-posterior}
p(\mathcal{M} | \mathcal{D}) = \frac{p(\mathcal{D} | \mathcal{M}) \, p(\mathcal{M})}{p(\mathcal{D})} \propto p(\mathcal{D} | \mathcal{M}) \, p(\mathcal{M})
\end{equation}
Under a uniform model prior ($p(\mathcal{M})$ constant), optimal model selection corresponds to maximizing the \textbf{marginal likelihood} (or model evidence):
\begin{equation}
\label{eq:dm-marginal-likelihood-integral}
p(\mathcal{D} | \mathcal{M}) = \int_{\Theta} p(\mathcal{D} | \boldsymbol{\theta}, \mathcal{M}) \, p(\boldsymbol{\theta} | \mathcal{M}) \, d\boldsymbol{\theta}
\end{equation}
Because the integral~(\ref{eq:dm-marginal-likelihood-integral}) is analytically intractable in general, Schwarz applied the \textbf{Laplace approximation} as sample size $R \to \infty$. Expanding the log-likelihood function $\ell(\boldsymbol{\theta}) = \ln p(\mathcal{D} | \boldsymbol{\theta}, \mathcal{M})$ in a second-order Taylor series around the Maximum Likelihood Estimator (MLE) $\hat{\boldsymbol{\theta}}$:
\begin{equation}
\label{eq:dm-taylor-laplace}
\ell(\boldsymbol{\theta}) \approx \ell(\hat{\boldsymbol{\theta}}) - \frac{1}{2} (\boldsymbol{\theta} - \hat{\boldsymbol{\theta}})^T \mathbf{H} (\boldsymbol{\theta} - \hat{\boldsymbol{\theta}})
\end{equation}
where $\mathbf{H} = -\nabla^2 \ell(\hat{\boldsymbol{\theta}})$ is the observed Fisher information matrix. Substituting this Gaussian quadratic form into the integral and evaluating the asymptotic determinant for large $R$ gives:
\begin{equation}
\label{eq:dm-bic-derivation-limit}
\ln p(\mathcal{D} | \mathcal{M}) = \ell(\hat{\boldsymbol{\theta}}) - \frac{p}{2} \ln R + \mathcal{O}(1)
\end{equation}
Defining $\mathrm{BIC}(\mathcal{M})$ proportional to the log marginal likelihood yields the standard formulation:
\begin{equation}
\label{eq:dm-bic-definition}
\mathrm{BIC}(\mathcal{M}) = \ell(\hat{\boldsymbol{\theta}}) - \frac{p}{2} \ln R
\end{equation}
where:
\begin{itemize}
    \item $\ell(\hat{\boldsymbol{\theta}}) = \sum_{i=1}^R \ln p(\mathbf{x}_i | \hat{\boldsymbol{\theta}})$ is the maximized log-likelihood of the data under the model;
    \item $p$ is the total count of free estimated parameters;
    \item $R$ is the number of observations;
    \item $-\frac{p}{2} \ln R$ is \textbf{Schwarz's complexity penalty}, growing logarithmically with sample size and linearly with parametric degrees of freedom.
\end{itemize}

\begin{remark}[Sign Conventions in Statistical Literature]
Two sign conventions coexist in literature:
\begin{enumerate}
    \item \textbf{Regularized Likelihood Convention} (standard in data mining and Pelleg-Moore): $\mathrm{BIC} = \ell(\hat{\boldsymbol{\theta}}) - \frac{p}{2} \ln R$, which is \textbf{maximized} (higher is better).
    \item \textbf{Deviance Convention} (common in econometrics): $\mathrm{BIC}_{\text{dev}} = -2 \ell(\hat{\boldsymbol{\theta}}) + p \ln R$, which is \textbf{minimized}. Both conventions are strictly equivalent as $\mathrm{BIC}_{\text{dev}} = -2 \cdot \mathrm{BIC}$.
\end{enumerate}
\end{remark}

\subsection{Formalization of BIC for Spherical Clustering}

In the $X$-Means framework~\cite{pelleg2000}, each cluster $k \in \{1, \dots, K\}$ is modeled as an isotropic spherical Gaussian in $d$ dimensions with shared variance $\boldsymbol{\Sigma}_k = \sigma^2 \mathbf{I}_d$. The probability density function of point $\mathbf{x}$ generated by cluster $\mathcal{C}_k$ is:
\begin{equation}
\label{eq:dm-spherical-gaussian-pdf}
p(\mathbf{x} | \mathbf{c}_k, \sigma^2) = \frac{1}{(2\pi)^{d/2} \sigma^d} \exp\left( -\frac{\|\mathbf{x} - \mathbf{c}_k\|_2^2}{2\sigma^2} \right)
\end{equation}
The log-likelihood of $R$ points partitioned across $K$ clusters is given analytically by:
\begin{equation}
\label{eq:dm-loglik-spherical-kmeans}
\ell(\hat{\boldsymbol{\theta}}) = -\frac{R \cdot d}{2} \ln(2\pi) - \frac{R \cdot d}{2} \ln(\hat{\sigma}^2) - \frac{1}{2\hat{\sigma}^2} \sum_{k=1}^K \mathrm{SSE}(\mathcal{C}_k) + \sum_{k=1}^K R_k \ln\left(\frac{R_k}{R}\right)
\end{equation}
where $R_k = |\mathcal{C}_k|$, and the unbiased estimator of the pooled residual variance $\hat{\sigma}^2$ accounts for the degrees of freedom:
\begin{equation}
\label{eq:dm-sigma-sq-unbiased}
\hat{\sigma}^2 = \frac{1}{R - K} \frac{1}{d} \sum_{k=1}^K \mathrm{SSE}(\mathcal{C}_k)
\end{equation}
Substituting $\hat{\sigma}^2$ into~(\ref{eq:dm-loglik-spherical-kmeans}) simplifies the third term to exactly $-\frac{(R - K) \cdot d}{2}$.

\subsection{Rigorous Counting of Free Parameters}

The parameter budget $p_K$ for a model with $K$ clusters consists of:
\begin{itemize}
    \item $K \cdot d$ parameters for the spatial coordinates of the $K$ centroid vectors $\mathbf{c}_k \in \mathbb{R}^d$;
    \item $1$ parameter for the shared isotropic variance $\hat{\sigma}^2$;
    \item $K - 1$ free parameters for cluster prior weights $\pi_k = \frac{R_k}{R}$ (since $\sum_{k=1}^K \pi_k = 1$).
\end{itemize}
Thus, the total parameter count is:
\begin{equation}
\label{eq:dm-pk-spherical-count}
p_K = K \cdot d + 1 + (K - 1) = K(d + 1)
\end{equation}
When prior weights are assumed uniform or empirically fixed, literature often uses the reduced budget $p_K = K \cdot d + 1$.

\begin{table}[htbp]
    \centering
    \small
    \renewcommand{\arraystretch}{1.35}
    \begin{tabularx}{\textwidth}{>{\raggedright\arraybackslash}p{2.8cm} >{\centering\arraybackslash}p{2.9cm} >{\centering\arraybackslash}p{2.6cm} Y}
        \toprule
        \textbf{Covariance Model} & \textbf{Form of $\boldsymbol{\Sigma}_k$} & \textbf{Parameters $p$} & \textbf{Geometric Properties} \\
        \midrule
        \textbf{Shared Spherical} & $\sigma^2 \mathbf{I}_d$ & $K \cdot d + 1$ & Isotropic spheres of identical radius across all clusters. \\
        \addlinespace[3pt]
        \textbf{Distinct Spherical} & $\sigma_k^2 \mathbf{I}_d$ & $K(d + 1)$ & Spheres with heterogeneous volumes across clusters. \\
        \addlinespace[3pt]
        \textbf{Diagonal (Independent)} & $\operatorname{diag}(\sigma_{k1}^2, \dots, \sigma_{kd}^2)$ & $K \cdot 2d$ & Hyper-ellipsoids aligned with coordinate axes. \\
        \addlinespace[3pt]
        \textbf{Full (General)} & Arbitrary $\boldsymbol{\Sigma}_k \succ 0$ & $K \left(d + \frac{d(d+1)}{2}\right)$ & Ellipsoids with arbitrary orientation, rotation, and eccentricity. \\
        \bottomrule
    \end{tabularx}
    \caption{Comparison of free parameter counts $p$ as a function of Gaussian covariance parameterization for $K$ components in dimension $d$.}
    \label{tab:dm-covariance-parameters}
\end{table}

Table~\ref{tab:dm-covariance-parameters} highlights the quadratic explosion in complexity ($\mathcal{O}(K \cdot d^2)$) when transitioning from spherical to full covariance matrices, demonstrating why $X$-Means employs the spherical formulation for computational scalability and numerical stability in high dimensions.

% ====================================================================
% SECTION 9.5: MIXTURE MODELS AND EXPECTATION-MAXIMIZATION (EM)
% ====================================================================
\section[Mixture Models and the EM Algorithm]{Mixture Models and the Expectation-Maximization (EM) Algorithm}
\label{sec:dm-mixture-em}

While geometric partitioning methods assign points deterministically, the probabilistic generative paradigm founded on \textbf{Mixture Models} and the \textbf{Expectation-Maximization} (EM) algorithm~\cite{dempster1977} views data as realizations of a random variable generated by a convex combination of elementary probability distributions.

\subsection{Generative Formulation and Latent Variables}

Each observation $\mathbf{x}_i \in \mathbb{R}^d$ ($i = 1, \dots, N$) is assumed to arise from the following stochastic process:
\begin{enumerate}
    \item A latent categorical component $z_i \in \{1, \dots, K\}$ is drawn according to prior probabilities:
    \begin{equation}
    p(z_i = k) = \pi_k, \qquad \text{with } \sum_{k=1}^K \pi_k = 1 \quad \text{and} \quad \pi_k \ge 0
    \end{equation}
    The values $\pi_k$ are termed \textbf{mixing proportions} (or weights).
    \item Conditioned on component $k$, the observation $\mathbf{x}_i$ is sampled from distribution $p_k(\mathbf{x}_i | \boldsymbol{\theta}_k)$:
    \begin{equation}
    \mathbf{x}_i \sim p(\mathbf{x}_i | z_i = k, \boldsymbol{\theta}_k)
    \end{equation}
\end{enumerate}
The marginal probability density of observation $\mathbf{x}_i$ is obtained via the law of total probability:
\begin{equation}
\label{eq:dm-mixture-marginal-density}
p(\mathbf{x}_i | \boldsymbol{\Theta}) = \sum_{k=1}^K p(z_i = k) \, p(\mathbf{x}_i | z_i = k, \boldsymbol{\theta}_k) = \sum_{k=1}^K \pi_k \, p_k(\mathbf{x}_i | \boldsymbol{\theta}_k)
\end{equation}
where $\boldsymbol{\Theta} = \{\pi_1, \dots, \pi_K, \boldsymbol{\theta}_1, \dots, \boldsymbol{\theta}_K\}$ represents the full parameter set.

\subsection{Incomplete vs. Complete Log-Likelihood}

Given an i.i.d. dataset $\mathbf{X} = \{\mathbf{x}_1, \dots, \mathbf{x}_N\}$, the observed \textbf{incomplete log-likelihood} is:
\begin{equation}
\label{eq:dm-incomplete-loglik}
\ln p(\mathbf{X} | \boldsymbol{\Theta}) = \sum_{i=1}^N \ln \left( \sum_{k=1}^K \pi_k \, p_k(\mathbf{x}_i | \boldsymbol{\theta}_k) \right)
\end{equation}
The presence of the logarithm of a summation in~(\ref{eq:dm-incomplete-loglik}) couples all parameters, preventing closed-form solutions for $\nabla_{\boldsymbol{\Theta}} \ln p(\mathbf{X} | \boldsymbol{\Theta}) = \mathbf{0}$.

If the latent indicator variables $\mathbf{Z} = \{\mathbf{z}_1, \dots, \mathbf{z}_N\}$ were observed, using indicator variables $z_{ik} = \mathbb{I}(z_i = k)$, the joint \textbf{complete log-likelihood} of $(\mathbf{X}, \mathbf{Z})$ would take the decoupled form:
\begin{equation}
\label{eq:dm-complete-loglik}
\ln p(\mathbf{X}, \mathbf{Z} | \boldsymbol{\Theta}) = \sum_{i=1}^N \sum_{k=1}^K z_{ik} \left[ \ln \pi_k + \ln p_k(\mathbf{x}_i | \boldsymbol{\theta}_k) \right]
\end{equation}
where the logarithm acts directly on individual component densities, permitting closed-form independent optimization.

\subsection{Variational Foundation of the EM Algorithm}

The EM algorithm optimizes the incomplete log-likelihood by applying Jensen's inequality to establish a variational lower bound computed under the posterior distribution of latent variables.

Let $\boldsymbol{\Theta}^{(t)}$ denote parameter estimates at step $t$. The auxiliary function $\mathcal{Q}(\boldsymbol{\Theta}, \boldsymbol{\Theta}^{(t)})$ is defined as the conditional expectation of the complete log-likelihood:
\begin{align}
\label{eq:dm-q-function}
\mathcal{Q}(\boldsymbol{\Theta}, \boldsymbol{\Theta}^{(t)}) &= \mathbb{E}_{\mathbf{Z} | \mathbf{X}, \boldsymbol{\Theta}^{(t)}} \left[ \ln p(\mathbf{X}, \mathbf{Z} | \boldsymbol{\Theta}) \right] \nonumber \\
&= \sum_{i=1}^N \sum_{k=1}^K \mathbb{E}[z_{ik} | \mathbf{x}_i, \boldsymbol{\Theta}^{(t)}] \left[ \ln \pi_k + \ln p_k(\mathbf{x}_i | \boldsymbol{\theta}_k) \right] \nonumber \\
&= \sum_{i=1}^N \sum_{k=1}^K \gamma_{ik} \left[ \ln \pi_k + \ln p_k(\mathbf{x}_i | \boldsymbol{\theta}_k) \right]
\end{align}
where $\gamma_{ik} \equiv \mathbb{E}[z_{ik} | \mathbf{x}_i, \boldsymbol{\Theta}^{(t)}] = p(z_i = k | \mathbf{x}_i, \boldsymbol{\Theta}^{(t)})$ represents the \textbf{posterior responsibility} that component $k$ takes for generating instance $\mathbf{x}_i$.

\begin{theorem}[Monotonicity of Likelihood in EM]
For any sequence of parameter vectors $\boldsymbol{\Theta}^{(t)}$ and $\boldsymbol{\Theta}^{(t+1)}$ generated by EM, the likelihood satisfies the monotonic ascent property:
\begin{equation}
\ln p(\mathbf{X} | \boldsymbol{\Theta}^{(t+1)}) \ge \ln p(\mathbf{X} | \boldsymbol{\Theta}^{(t)})
\end{equation}
The algorithm converges monotonically to a stationary point (local maximum or saddle point) of the incomplete log-likelihood landscape.
\end{theorem}

\subsection{Operational Cycle: E-Step and M-Step}

EM iterates through two steps until numerical convergence:

\paragraph{1. Expectation Step (E-step)}
Holding parameters $\boldsymbol{\Theta}^{(t)}$ fixed, apply Bayes' rule to compute posterior responsibilities:
\begin{equation}
\label{eq:dm-e-step}
\gamma_{ik} = \frac{\pi_k^{(t)} \, p_k(\mathbf{x}_i | \boldsymbol{\theta}_k^{(t)})}{\sum_{j=1}^K \pi_j^{(t)} \, p_j(\mathbf{x}_i | \boldsymbol{\theta}_j^{(t)})}
\end{equation}
Note that $\sum_{k=1}^K \gamma_{ik} = 1$ for every observation $i$. The \textbf{effective sample size} belonging to cluster $k$ is:
\begin{equation}
\label{eq:dm-effective-nk}
N_k = \sum_{i=1}^N \gamma_{ik}
\end{equation}

\paragraph{2. Maximization Step (M-step)}
Holding responsibilities $\gamma_{ik}$ fixed, obtain updated parameters $\boldsymbol{\Theta}^{(t+1)}$ by maximizing $\mathcal{Q}$:
\begin{equation}
\label{eq:dm-m-step-argmax}
\boldsymbol{\Theta}^{(t+1)} = \arg\max_{\boldsymbol{\Theta}} \mathcal{Q}(\boldsymbol{\Theta}, \boldsymbol{\Theta}^{(t)})
\end{equation}
For Gaussian Mixture Models (GMM), closed-form updates for mixing weights, means, and covariances are:
\begin{align}
\label{eq:dm-m-step-pi}
\pi_k^{(t+1)} &= \frac{N_k}{N} = \frac{1}{N} \sum_{i=1}^N \gamma_{ik} \\
\label{eq:dm-m-step-mu}
\boldsymbol{\mu}_k^{(t+1)} &= \frac{1}{N_k} \sum_{i=1}^N \gamma_{ik} \mathbf{x}_i \\
\label{eq:dm-m-step-sigma}
\boldsymbol{\Sigma}_k^{(t+1)} &= \frac{1}{N_k} \sum_{i=1}^N \gamma_{ik} (\mathbf{x}_i - \boldsymbol{\mu}_k^{(t+1)}) (\mathbf{x}_i - \boldsymbol{\mu}_k^{(t+1)})^T
\end{align}

\begin{figure}[htbp]
\centering
\begin{tikzpicture}[
    box/.style={rectangle, rounded corners=4.5pt, draw=navy!80!black, fill=navy!4, thick, font=\footnotesize, inner sep=6pt, align=left, text width=11.0cm},
    initBox/.style={rectangle, rounded corners=4.5pt, draw=navy!70!black, fill=navy!3, thick, font=\footnotesize, inner sep=5pt, align=left, text width=11.0cm},
    checkCard/.style={rectangle, rounded corners=4.5pt, draw=orange!85!black, fill=orange!8, thick, font=\footnotesize, inner sep=5pt, align=left, text width=11.0cm},
    resultBadge/.style={rectangle, rounded corners=4pt, draw=darkgreen!85!black, fill=darkgreen!10, thick, font=\footnotesize\bfseries, text=darkgreen!90!black, inner sep=5pt, align=center, text width=9.2cm},
    tagBadge/.style={rectangle, rounded corners=3.5pt, draw=purple!85!black, fill=purple!12, thick, font=\scriptsize\bfseries, text=purple!90!black, inner sep=2.5pt},
    flowArrow/.style={-Stealth, very thick, draw=navy!80!black},
    loopArrow/.style={-Stealth, very thick, draw=crimson!85!black}
]
    % Node 1: Initialization
    \node[initBox] (init) at (0, 0) {
        \textbf{Parameter Initialization $\boldsymbol{\Theta}^{(0)}$}\\[2pt]
        $\bullet$ Uniform mixing weights $\pi_k^{(0)} = \frac{1}{K}$, initial means $\boldsymbol{\mu}_k^{(0)}$ (from $K$-Means or random).\\
        $\bullet$ Initial isotropic covariance matrices: $\boldsymbol{\Sigma}_k^{(0)} = \hat{\sigma}^2 \mathbf{I}_d$.
    };

    % Node 2: E-Step
    \node[box, below = 0.55cm of init] (estep) {
        \textbf{E-Step (Expectation)}\\[2pt]
        $\bullet$ Compute posterior responsibilities and effective sample size given parameters $\boldsymbol{\Theta}^{(t)}$:
        \vspace{2pt}
        $$\gamma_{ik} = \frac{\pi_k^{(t)} p(\mathbf{x}_i \mid \boldsymbol{\theta}_k^{(t)})}{\sum_{j=1}^K \pi_j^{(t)} p(\mathbf{x}_i \mid \boldsymbol{\theta}_j^{(t)})}, \qquad N_k = \sum_{i=1}^N \gamma_{ik}$$
    };
    \node[tagBadge, anchor=south east] at ([xshift=-10pt, yshift=-1.5pt]estep.north east) {Stage 1};

    % Node 3: M-Step
    \node[box, below = 0.55cm of estep] (mstep) {
        \textbf{M-Step (Maximization)}\\[2pt]
        $\bullet$ Update parameters $\boldsymbol{\Theta}^{(t+1)}$ by maximizing expectation auxiliary function $\mathcal{Q}(\boldsymbol{\Theta}, \boldsymbol{\Theta}^{(t)})$:
        \vspace{2pt}
        $$\pi_k^{(t+1)} = \frac{N_k}{N}, \qquad \boldsymbol{\mu}_k^{(t+1)} = \frac{1}{N_k}\sum_{i=1}^N \gamma_{ik}\mathbf{x}_i$$
        \vspace{-6pt}
        $$\boldsymbol{\Sigma}_k^{(t+1)} = \frac{1}{N_k}\sum_{i=1}^N \gamma_{ik} (\mathbf{x}_i - \boldsymbol{\mu}_k^{(t+1)})(\mathbf{x}_i - \boldsymbol{\mu}_k^{(t+1)})^T$$
    };
    \node[tagBadge, anchor=south east] at ([xshift=-10pt, yshift=-1.5pt]mstep.north east) {Stage 2};

    % Node 4: Convergence Check
    \node[checkCard, below = 0.55cm of mstep] (check) {
        \textbf{Convergence Evaluation}\\[2pt]
        $\bullet$ Incomplete log-likelihood gain: $|\ln p(\mathbf{X}\mid\boldsymbol{\Theta}^{(t+1)}) - \ln p(\mathbf{X}\mid\boldsymbol{\Theta}^{(t)})| < \varepsilon$.
    };

    % Node 5: Optimal Model
    \node[resultBadge, below = 0.55cm of check] (result) {
        \textbf{Convergence Reached (Stationary Point)}\\
        \scriptsize Optimal Model $\boldsymbol{\Theta}^* = \{\pi_k^*, \boldsymbol{\mu}_k^*, \boldsymbol{\Sigma}_k^*\}$
    };

    % Forward Arrows
    \draw[flowArrow] (init.south) -- (estep.north)
        node[midway, right=4pt, font=\scriptsize\bfseries, text=navy!85!black] {$\boldsymbol{\Theta}^{(0)}$};

    \draw[flowArrow] (estep.south) -- (mstep.north)
        node[midway, right=4pt, font=\scriptsize\bfseries, text=navy!85!black] {Responsibilities $\gamma_{ik}$};

    \draw[flowArrow] (mstep.south) -- (check.north)
        node[midway, right=4pt, font=\scriptsize\bfseries, text=navy!85!black] {$\boldsymbol{\Theta}^{(t+1)}$};

    \draw[-Stealth, very thick, draw=darkgreen!85!black] (check.south) -- (result.north)
        node[midway, right=4pt, font=\scriptsize\bfseries, text=darkgreen!85!black] {Yes (stationary)};

    % Feedback Loop Arrow (on the left)
    \draw[loopArrow] (check.west) -- ++(-1.1, 0)
        node[midway, above=2pt, font=\scriptsize\bfseries, text=crimson!85!black] {No}
        |- ([yshift=-2pt]estep.west)
        node[pos=0.35, left=3pt, font=\scriptsize\bfseries, text=crimson!85!black, align=right] {New iteration $t \leftarrow t+1$\\with $\boldsymbol{\Theta}^{(t+1)}$};

\end{tikzpicture}
\caption{Operational flowchart of the Expectation-Maximization (EM) algorithm.}
\label{fig:dm-em-cycle}
\end{figure}

% ====================================================================
% SECTION 9.6: NORMAL DISTRIBUTIONS AND MLE ESTIMATION
% ====================================================================
\section{Normal Distributions and Statistical Estimation Foundations}
\label{sec:dm-gaussian-mle}

To fully contextualize the M-step updates in GMMs, we review Maximum Likelihood Estimation (MLE) for a univariate Gaussian.

The probability density function for a univariate normal random variable $x \sim \mathcal{N}(\mu, \sigma^2)$ is:
\begin{equation}
\label{eq:dm-1d-gaussian-pdf}
p(x | \mu, \sigma^2) = \frac{1}{\sqrt{2\pi}\sigma} \exp\left( -\frac{(x - \mu)^2}{2\sigma^2} \right)
\end{equation}
Given an i.i.d. sample of $N$ scalar observations $\{x_1, \dots, x_N\}$, the log-likelihood function is:
\begin{equation}
\label{eq:dm-1d-gaussian-loglik}
\ell(\mu, \sigma^2) = \sum_{i=1}^N \ln p(x_i | \mu, \sigma^2) = -\frac{N}{2} \ln(2\pi) - \frac{N}{2} \ln(\sigma^2) - \frac{1}{2\sigma^2} \sum_{i=1}^N (x_i - \mu)^2
\end{equation}
Setting the partial derivative with respect to $\mu$ to zero:
\begin{equation}
\frac{\partial \ell}{\partial \mu} = \frac{1}{\sigma^2} \sum_{i=1}^N (x_i - \mu) = 0 \implies \hat{\mu}_{\text{MLE}} = \frac{1}{N} \sum_{i=1}^N x_i
\end{equation}
yields the sample arithmetic mean. Substituting $\hat{\mu}$ and setting the partial derivative with respect to $\sigma^2$ to zero:
\begin{equation}
\frac{\partial \ell}{\partial (\sigma^2)} = -\frac{N}{2\sigma^2} + \frac{1}{2(\sigma^2)^2} \sum_{i=1}^N (x_i - \hat{\mu})^2 = 0 \implies \hat{\sigma}^2_{\text{MLE}} = \frac{1}{N} \sum_{i=1}^N (x_i - \hat{\mu})^2
\end{equation}
In GMMs, the M-step extends these derivations directly by substituting simple arithmetic averages with \textbf{weighted averages} governed by posterior responsibilities $\gamma_{ik}$ (Equations~(\ref{eq:dm-m-step-mu}) and~(\ref{eq:dm-m-step-sigma})).

\paragraph{Bayesian MAP vs. MLE Foundations}
While MLE maximizes only the data likelihood $\hat{\boldsymbol{\Theta}}_{\text{MLE}} = \arg\max \ln p(\mathbf{X}|\boldsymbol{\Theta})$, \textbf{Maximum A Posteriori (MAP)} estimation incorporates conjugate parameter priors $p(\boldsymbol{\Theta})$:
\begin{equation}
\hat{\boldsymbol{\Theta}}_{\text{MAP}} = \arg\max_{\boldsymbol{\Theta}} \big[ \ln p(\mathbf{X}|\boldsymbol{\Theta}) + \ln p(\boldsymbol{\Theta}) \big]
\end{equation}
In high-dimensional mixture modeling, MAP estimation with Normal-Inverse-Wishart priors prevents singular covariance determinants ($|\boldsymbol{\Sigma}_k| \to 0$), acting as structural regularization (\textit{ridge shrinkage}).

% ====================================================================
% SECTION 9.7: CASE STUDY: TWO-GAUSSIAN MIXTURE
% ====================================================================
\section{Case Study: Deconvolving a Two-Gaussian Mixture}
\label{sec:dm-case-study-two-gaussians}

To gain concrete intuition for probabilistic soft clustering, consider the motivating educational scenario discussed in~\cite{tan2018}: modeling human height distributions in a mixed demographic population (e.g., overlapping subpopulations with distinct average stature).

\subsection{Definition of the Bimodal Mixture Model}

Consider a balanced univariate mixture ($K=2$) of two normal distributions with identical variance but shifted centers of mass:
\begin{itemize}
    \item \textbf{Component 1} ($\mathcal{C}_1$): Mean $\mu_1 = -4$, standard deviation $\sigma_1 = 2$, mixing weight $\pi_1 = 0.5$;
    \item \textbf{Component 2} ($\mathcal{C}_2$): Mean $\mu_2 = +4$, standard deviation $\sigma_2 = 2$, mixing weight $\pi_2 = 0.5$.
\end{itemize}
The marginal probability density of the combined population is:
\begin{align}
\label{eq:dm-two-gaussian-pdf}
p(x) &= 0.5 \cdot \mathcal{N}(x \,|\, -4, 4) + 0.5 \cdot \mathcal{N}(x \,|\, +4, 4) \nonumber \\
&= \frac{1}{2\sqrt{8\pi}} \left[ \exp\left( -\frac{(x+4)^2}{8} \right) + \exp\left( -\frac{(x-4)^2}{8} \right) \right]
\end{align}

\begin{figure}[htbp]
\centering
\begin{tikzpicture}[
    x=0.75cm, y=32cm,
    declare function={
        gauss1(\x) = 0.5 * (1.0 / (2.0 * sqrt(2.0 * 3.1415926))) * exp(-((\x + 4.0)^2) / 8.0);
        gauss2(\x) = 0.5 * (1.0 / (2.0 * sqrt(2.0 * 3.1415926))) * exp(-((\x - 4.0)^2) / 8.0);
        gaussTot(\x) = gauss1(\x) + gauss2(\x);
    }
]
    % Subtle horizontal grid lines
    \foreach \yval in {0.02, 0.04, 0.06, 0.08, 0.10} {
        \draw[dotted, gray!40] (-8.5, \yval) -- (8.5, \yval);
        \node[left=3pt, font=\scriptsize, text=gray!80!black] at (-8.5, \yval) {$\yval$};
    }

    % Shaded areas for each component (translucent, showing overlap)
    \fill[navy!15, opacity=0.6, domain=-8.5:3.0, samples=80]
        (-8.5, 0) -- plot (\x, {gauss1(\x)}) -- (3.0, 0) -- cycle;
    \fill[crimson!15, opacity=0.6, domain=-3.0:8.5, samples=80]
        (-3.0, 0) -- plot (\x, {gauss2(\x)}) -- (8.5, 0) -- cycle;

    % Component 1 curve (dashed blue)
    \draw[thick, dashed, draw=navy!85!blue, domain=-8.5:3.0, samples=80, smooth]
        plot (\x, {gauss1(\x)});

    % Component 2 curve (dashed red)
    \draw[thick, dashed, draw=crimson!85!black, domain=-3.0:8.5, samples=80, smooth]
        plot (\x, {gauss2(\x)});

    % Total mixture curve (solid, thick purple)
    \draw[very thick, draw=purple!90!black, domain=-8.5:8.5, samples=120, smooth]
        plot (\x, {gaussTot(\x)});

    % Vertical drop lines to means
    \draw[dashed, navy!70!black, thick] (-4, 0) -- (-4, {gaussTot(-4)});
    \draw[dashed, crimson!70!black, thick] (4, 0) -- (4, {gaussTot(4)});

    % Peak 1 point & label
    \draw[fill=navy!85!blue, draw=white, thick] (-4, {gaussTot(-4)}) circle (2.5pt);
    \node[above=3pt, font=\scriptsize\bfseries, text=navy!85!black, fill=white, inner sep=1.5pt, rounded corners=1.5pt]
        at (-4, {gaussTot(-4)}) {Mode $\mathcal{C}_1$ ($\mu_1 = -4$)};

    % Peak 2 point & label
    \draw[fill=crimson!85!black, draw=white, thick] (4, {gaussTot(4)}) circle (2.5pt);
    \node[above=3pt, font=\scriptsize\bfseries, text=crimson!85!black, fill=white, inner sep=1.5pt, rounded corners=1.5pt]
        at (4, {gaussTot(4)}) {Mode $\mathcal{C}_2$ ($\mu_2 = +4$)};

    % Decision boundary / Critical overlap line at x = 0
    \draw[dotted, thick, draw=purple!85!black] (0, 0) -- (0, 0.055);
    \draw[fill=purple!85!black, draw=white, thick] (0, {gaussTot(0)}) circle (2.5pt);

    % Callout for critical overlap
    \node[draw=purple!85!black, fill=white, rounded corners=3pt, inner sep=3.5pt, align=center, font=\scriptsize]
        (overlapNode) at (0, 0.075) {
            \textbf{Decision Boundary ($x = 0$)}\\[1.5pt]
            $\gamma_{1}(0) = \gamma_{2}(0) = 0.50$ (Bayesian Indecision)
        };
    \draw[-Stealth, thick, draw=purple!85!black] (overlapNode.south) -- (0, 0.033);

    % Coordinate axes
    % y-axis on the left
    \draw[-Stealth, thick, draw=navy!90!black] (-8.5, 0) -- (-8.5, 0.125)
        node[above=3pt, font=\footnotesize\bfseries, text=navy!90!black] {$p(x)$};
    % x-axis along bottom
    \draw[-Stealth, thick, draw=navy!90!black] (-8.5, 0) -- (8.8, 0)
        node[right=3pt, font=\footnotesize\bfseries, text=navy!90!black] {$x$};

    % x-axis ticks
    \foreach \xval in {-8, -6, -4, -2, 0, 2, 4, 6, 8} {
        \draw[thick, draw=navy!90!black] (\xval, 0) -- (\xval, -0.003);
        \node[below=3pt, font=\scriptsize, text=gray!90!black] at (\xval, 0) {$\xval$};
    }

    % Decision regions at bottom
    \node[below=14pt, font=\scriptsize, text=navy!80!black] at (-4, 0) {$\longleftarrow$ Decision Region $\mathcal{C}_1$ ($\gamma_1 > 0.5$)};
    \node[below=14pt, font=\scriptsize, text=crimson!80!black] at (4, 0) {Decision Region $\mathcal{C}_2$ ($\gamma_2 > 0.5$) $\longrightarrow$};

    % Legend box (Top center, above peaks)
    \node[draw=gray!50, fill=white, rounded corners=3.5pt, font=\scriptsize, inner sep=4pt] at (0, 0.125) {
        \textcolor{purple!90!black}{\rule{14pt}{2.0pt}} Total Mixture $p(x)$ \qquad
        \textcolor{navy!85!blue}{\rule{14pt}{1.3pt}} Comp. $\mathcal{C}_1$ ($\mu_1 = -4, \sigma = 2$) \qquad
        \textcolor{crimson!85!black}{\rule{14pt}{1.3pt}} Comp. $\mathcal{C}_2$ ($\mu_2 = +4, \sigma = 2$)
    };
\end{tikzpicture}
\caption{Analytical deconvolution of a bimodal two-Gaussian mixture with equal variance ($\sigma=2$) and separated centers ($\Delta\mu = 8$).}
\label{fig:dm-two-gaussian-density}
\end{figure}

\subsection{Overlap Analysis and Posterior Responsibilities}

Figure~\ref{fig:dm-two-gaussian-density} illustrates the key mathematical properties:
\begin{enumerate}
    \item \textbf{Two Distinct Local Maxima}: Clear peaks appear near $x = -4$ and $x = +4$.
    \item \textbf{Saddle Point and Overlap Uncertainty at $x = 0$}: At the midpoint $x = 0$, both conditional densities are identical:
    \begin{equation}
    p_1(0) = \frac{1}{\sqrt{8\pi}} \exp\left(-\frac{16}{8}\right) = \frac{e^{-2}}{\sqrt{8\pi}}, \qquad p_2(0) = \frac{e^{-2}}{\sqrt{8\pi}}
    \end{equation}
    Consequently, by Equation~(\ref{eq:dm-e-step}), posterior responsibilities are equal:
    \begin{equation}
    \gamma_{1}(0) = \frac{0.5 \cdot p_1(0)}{0.5 \cdot p_1(0) + 0.5 \cdot p_2(0)} = 0.5, \qquad \gamma_{2}(0) = 0.5
    \end{equation}
    At this location, any rigid deterministic cutoff (assigning $\mathcal{C}_1$ if $x < 0$ and $\mathcal{C}_2$ if $x \ge 0$) artificially conceals epistemic ambiguity. EM transparently reflects this uncertainty through continuous responsibilities $\gamma_{ik} = 0.5$.
    \item \textbf{Validation on 20,000 Empirical Samples}: When $N = 20\,000$ synthetic points are drawn from mixture~(\ref{eq:dm-two-gaussian-pdf}), the empirical histogram matches the continuous theoretical curve within Central Limit Theorem bounds. EM initialized with perturbed values accurately recovers true generating parameters $(\mu_1, \mu_2, \sigma_1, \sigma_2, \pi_1, \pi_2)$ with minimal parameter estimation error ($< 10^{-3}$).
\end{enumerate}

% ====================================================================
% SECTION 9.8: ARCHITECTURAL COMPARISON: K-MEANS VS. EM
% ====================================================================
\section{Architectural and Systematic Comparison: \texorpdfstring{$K$}{K}-Means vs. EM}
\label{sec:dm-kmeans-em-comparison}

Lloyd's algorithm ($K$-Means) and the EM algorithm for Gaussian Mixture Models share a formally isomorphic iterative architecture: both start with parameter estimates and alternate between updating point cluster memberships and recalculating cluster parameters. However, they diverge fundamentally in their geometric, energetic, and statistical foundations.

\begin{table}[htbp]
    \centering
    \small
    \renewcommand{\arraystretch}{1.3}
    \begin{tabularx}{\textwidth}{>{\raggedright\arraybackslash\bfseries}p{3.3cm} Y Y}
        \toprule
        \textbf{Dimension} & \textbf{Standard $K$-Means (Lloyd)} & \textbf{Gaussian Mixture (EM)} \\
        \midrule
        Assignment Paradigm & \textbf{Hard clustering}: deterministic and disjoint ($r_{ik} \in \{0, 1\}$). & \textbf{Soft clustering}: probabilistic and continuous ($\gamma_{ik} \in [0, 1]$). \\
        \addlinespace[3pt]
        Cluster Model & Barycentric centroid vector $\mathbf{c}_k \in \mathbb{R}^d$. & Parameter tuple $(\pi_k, \boldsymbol{\mu}_k, \boldsymbol{\Sigma}_k)$. \\
        \addlinespace[3pt]
        Geometric Shape & Hyperspherical, isotropic, homogeneous variance. & Hyper-ellipsoidal with arbitrary orientation, size, and shape. \\
        \addlinespace[3pt]
        Objective Function & Minimizing intra-cluster Sum of Squared Errors ($\mathrm{SSE}$). & Maximizing incomplete log-likelihood ($\ln p(\mathbf{X}|\boldsymbol{\Theta})$). \\
        \addlinespace[3pt]
        Phase 1 (Point State) & Voronoi assignment to nearest centroid ($L_2$). & E-step: analytical posterior responsibilities $\gamma_{ik}$ (Bayes). \\
        \addlinespace[3pt]
        Phase 2 (Update) & Unweighted arithmetic mean of assigned points. & M-step: weighted mean and covariance using weights $\gamma_{ik}$. \\
        \addlinespace[3pt]
        Computational Cost & $\mathcal{O}(N \cdot K \cdot d)$ & $\mathcal{O}(N \cdot K \cdot d^2 + K \cdot d^3)$ (matrix inversion). \\
        \addlinespace[3pt]
        Vulnerabilities & Seed sensitivity, spherical bias, zero uncertainty quantification. & Covariance singularities, $\mathcal{O}(d^3)$ cost, local optima. \\
        \bottomrule
    \end{tabularx}
    \caption{Comprehensive architectural comparison between $K$-Means and Expectation-Maximization (EM).}
    \label{tab:dm-kmeans-em-detailed-comparison}
\end{table}

\subsection{Covariance Singularities in EM}

A significant vulnerability of EM with unconstrained covariance matrices $\boldsymbol{\Sigma}_k$ is the existence of \textbf{singularities in the likelihood surface}. If a Gaussian component collapses onto an isolated singleton point $\mathbf{x}_m$, its sample covariance shrinks to zero:
\begin{equation}
\boldsymbol{\Sigma}_k \to \mathbf{0} \implies |\boldsymbol{\Sigma}_k| \to 0
\end{equation}
Because the multivariate Gaussian density includes $|\boldsymbol{\Sigma}_k|^{1/2}$ in its denominator:
\begin{equation}
p_k(\mathbf{x}_m | \boldsymbol{\mu}_k, \boldsymbol{\Sigma}_k) \to +\infty \implies \ln p(\mathbf{X} | \boldsymbol{\Theta}) \to +\infty
\end{equation}
The likelihood function is unbounded from above, generating infinite spikes on the surface that represent pathological overfitting to single observations rather than meaningful clusters. Common remedies include:
\begin{itemize}
    \item \textbf{Diagonal Jitter}: Adding a small positive constant $\varepsilon \mathbf{I}_d$ to covariance matrices at every step ($\boldsymbol{\Sigma}_k' = \boldsymbol{\Sigma}_k + \epsilon \mathbf{I}_d$).
    \item \textbf{MAP Regularization}: Imposing an Inverse-Wishart prior on $\boldsymbol{\Sigma}_k$ to penalize shrinking determinants.
\end{itemize}

\subsection{Summary of the Methodological Path}

The conceptual journey presented in this chapter charts the modern evolution of clustering:
\begin{itemize}
    \item Divisive hierarchical clustering (\textbf{Bisecting $K$-Means}) converts an unstable global heuristic into a stable sequence of low-complexity binary decisions.
    \item Statistical information theory (\textbf{$X$-Means} and \textbf{BIC}) replaces subjective guesses for $K$ with objective complexity optimization rooted in Occam's razor.
    \item Probabilistic generative modeling (\textbf{GMM} and \textbf{EM}) elevates spatial partitioning into rigorous statistical inference, quantifying ambiguity and modeling arbitrary oriented geometries in continuous space.
\end{itemize}
